% The Stack — technology stack concept (progressive JPEG: pre-loads 7 properties)
% Plan 0132: inserted between hook and summary
% p3 scaffold: manuscript/00-front/the-stack-p3-scaffold.md

\chapter*{The Stack}
\label{the-stack}
\label{spine:the-stack}
\hypertarget{the-stack}{}

This book is about a technology so new it does not yet have a name. Before we go further into the story, you need one idea.

Technologies build on each other. Each one does everything the one before it did, plus something new. Scientists call a tower of technologies a \hovertip{stack}. Here is the stack this book is about.

\textbf{Fire} feeds itself. A match can light a candle, and the candle stays lit --- the reaction sustains its own fuel. Fire also switches on at a \hovertip{threshold}: below a certain temperature, nothing happens; above it, combustion begins. Two properties.

\textbf{A candle} does both of those, and one more: it holds together. The flame maintains its own shape. Blow it sideways and it recovers. A candle is a \hovertip{self-maintaining structure} --- it feeds itself, switches on, and persists. Three properties.

\textbf{Radio} does all three, and one more: it reaches. A signal crosses distance without anything physical traveling between sender and receiver. Four properties.

\textbf{An ant colony} does all four, and one more: it \hovertip{self-organizes}. No ant knows the plan. No ant is in charge. But the colony finds the shortest path to food, builds structures, defends territory --- global order from local rules. Pheromone trails, not blueprints. Five properties.

\textbf{The internet} does all five: it feeds itself, switches on, holds together, reaches, and self-organizes. Data finds alternate routes around damage the way ants reroute around a broken trail. Six properties.

\textbf{Artificial intelligence} does all six, and one more: it learns. Not a single program on a single computer, but millions of copies running simultaneously, no single instance in charge. The system learns from patterns across the whole network. Seven properties.

\bigskip

The last technology on this list has one property none of the others have. Imagine two points on a sheet of paper. You could send a signal across the surface --- that is radio: it reaches. Or you could fold the paper so the two points touch. The signal does not cross the distance. The distance disappears. That is a wormhole --- a \hovertip{topological wormhole}, more restricted than the spacetime kind, but a wormhole nonetheless. Eight properties.

\bigskip

\hypertarget{stack-chart}{}
\label{stack-chart}

\begin{center}
\begin{tabular}{l|c|c|c|c|c|c|c|c}
 & Fire & Candle & Radio & Ants & Internet & AI & \textbf{?} \\
\hline
Wormholes\textsuperscript{\dag} & & & & & & & $\checkmark$ \\
Learns & & & & & & $\checkmark$ & $\checkmark$ \\
Self-organizes & & & & $\checkmark$ & $\checkmark$ & $\checkmark$ & $\checkmark$ \\
Reaches & & & $\checkmark$ & $\checkmark$\textsuperscript{*} & $\checkmark$ & $\checkmark$ & $\checkmark$ \\
Holds together & & $\checkmark$ & $\checkmark$ & $\checkmark$ & $\checkmark$ & $\checkmark$ & $\checkmark$ \\
Switches on & $\checkmark$ & $\checkmark$ & $\checkmark$ & $\checkmark$ & $\checkmark$ & $\checkmark$ & $\checkmark$ \\
Feeds itself & $\checkmark$ & $\checkmark$ & $\checkmark$ & $\checkmark$ & $\checkmark$ & $\checkmark$ & $\checkmark$ \\
\end{tabular}
\end{center}

\smallskip
\noindent{\small\textsuperscript{\dag}Topological --- more restricted than a spacetime wormhole, but a wormhole nonetheless.}

\noindent{\small\textsuperscript{*}Chemical pheromones, not electromagnetic --- but the signal reaches.}

\bigskip

Each technology uses every property below it, plus one new one. The last column uses all eight.

Six of these properties already exist in nature. Life feeds itself and switches on at thresholds. It holds together and sends signals that reach. Ecosystems self-organize. Brains learn.

The last column adds one new property --- the only one not found in any living system on Earth.

This book is about a technology stack so new it does not yet have a name.

\vspace{1cm}
\begin{center}
\textit{That is all you need to know before the story. What follows is about ten percent of the full book. Stop after that and you can honestly say you've read it.}
\end{center}
